\begin{recipe}
[% 
    preparationtime = {\unit[10]{min}},
    bakingtime = {\unit[40]{min}},
    bakingtemperature = {\protect\bakingtemperature{
        topbottomheat=\unit[180]{\textcelcius}
    }},
    portion = {\portion{12}},
    source = {\protect\recipesource{https://biancazapatka.com/de/saftiger-und-einfacher-karottenkuchen-moehrenkuchen-rueblikuchen-vegan/}{Bianca Zapatka: Karottenkuchen}},
    calory = {285kcal | 4g Protein | 30g KH | 16g Fett},
    price = {0,78€},
    created = {05.05.2026},
    rating = {2},
]
{Veganer Karottenkuchen}

\graph{% pictures
    big=pic/dessert/11_veganer-karottenkuchen
}

\introduction{%
Saftiger veganer Karottenkuchen mit gemahlenen Mandeln und Frischkäse-Frosting. Eigentlich ein Osterklassiker, aber ganz ehrlich: Karotten gibt es das ganze Jahr, also gibt es auch keinen Grund, den nur einmal im Jahr zu backen.
}

\ingredients{
    \textbf{Teig} & \\
    250 g & Karotten\index{Gemüse \& Pilze!Karotten}\\
    125 g & Vegane Butter\index{Milchprodukte \& Alternativen!Vegane Butter}\\
    100 g & Zucker\index{Backzutaten \& Süßes!Zucker}\\
    1 Prise & Salz\index{Kräuter \& Gewürze!Salz}\\
    1 Prise & Gemahlene Vanille\index{Backzutaten \& Süßes!Vanille}\\
    2 & Leinsamen-Eier\index{Eier!Leinsamen-Eier}\index{Nüsse \& Samen!Leinsamen}\\
    \nicefrac{1}{2} TL & Zimt\index{Kräuter \& Gewürze!Zimt}\\
    100 g & Gemahlene Mandeln\index{Nüsse \& Samen!Mandeln}\\
    125 g & Weizenmehl\index{Getreide, Pasta \& Brot!Weizenmehl}\\
    1 TL & Apfelessig\index{Öle, Saucen \& Würzmittel!Apfelessig}\\
    1{,}5 TL & Backpulver\index{Backzutaten \& Süßes!Backpulver}\\
    \nicefrac{1}{2} TL & Natron\index{Backzutaten \& Süßes!Natron}\\
    \textbf{Creme} & \\
    200 g & Veganer Frischkäse\index{Milchprodukte \& Alternativen!Frischkäse}\\
    80 g & Puderzucker\index{Backzutaten \& Süßes!Puderzucker}\\
    1 Spritzer & Zitronensaft\index{Obst!Zitronen}\\
    Etwas & Gemahlene Vanille\index{Backzutaten \& Süßes!Vanille}\\
    \textbf{Optional} & \\
    & Gehackte Pistazien\index{Nüsse \& Samen!Pistazien}\\
    & Marzipan-Karotten\index{Backzutaten \& Süßes!Marzipan}\index{Gemüse \& Pilze!Karotten}
}

\preparation{%
    \step%
    Backofen auf 180 °C Ober-/Unterhitze vorheizen. Eine 18--20 cm Springform einfetten oder mit Backpapier auslegen.
    
    \step%
    Für die Leinsamen-Eier 2 EL gemahlene Leinsamen mit 6 EL heißem Wasser verrühren und 5--10 Minuten quellen lassen.
    
    \step%
    Karotten schälen und fein reiben.
    
    \step%
    Vegane Butter mit Zucker, Salz und Vanille cremig aufschlagen. Karotten und Leinsamen-Eier unterrühren.
    
    \step%
    Mehl, Zimt, Backpulver, gemahlene Mandeln, Natron und Apfelessig zugeben und nur kurz zu einem Teig verrühren.
    
    \step%
    Teig in die Form füllen, glatt streichen und 35--40 Minuten backen. Stäbchenprobe machen.
    
    \step%
    Kuchen 10 Minuten in der Form abkühlen lassen, dann herauslösen und vollständig auskühlen lassen.
    
    \step%
    Für die Creme veganen Frischkäse, Puderzucker, Zitronensaft und Vanille glatt rühren.
    
    \step%
    Creme auf den vollständig abgekühlten Kuchen streichen und nach Belieben mit Pistazien oder Marzipan-Karotten dekorieren.
}

\hint{
    Den Teig nach dem Zugeben von Mehl nur kurz rühren, sonst wird der Kuchen kompakter. Und ja, vollständig abkühlen lassen vor der Creme ist leider wirklich wichtig.
}

\suggestion[Alternative]{%
    Statt Leinsamen-Eiern funktionieren laut Original auch 100 g Apfelmark, Bananenpüree oder Kürbispüree.
}

\end{recipe}